\subsection*{Writing from a blank answer strip}
We tested whether seeing both sides of a missing word would help a visual writer compose a reply. The model predicted an answer length from a pictured passage and question, then filled the blank strip in four passes. When asked to identify a senior researcher, it wrote the first name \emph{Amanda} on pass one and the surname \emph{Lenhart} on pass three (Fig.~\ref{fig:word-canvas}). Among 3,155 questions from 200 previously unused stories, it matched 693 human answers, falling to 63 after question-and-history exchange and 353 after passage exchange. Its unchanged parent matched 769, and a matched left-to-right continuation matched 743.

On 443 questions whose reference required multiword phrasing absent from the inputs, the masked visual writer returned zero exact answers. The left-to-right control returned one: it assembled \emph{he killed himself} from three word images. The answer lost credit after either exchange. Difficult-turn word F1 was 13.03\% for the masked writer and 15.54\% for the causal control (difference $-2.51$ percentage points, 95\% paired story interval $-4.52$ to $-0.44$). The masked byte writer matched one difficult reference. On 279 questions whose normalized reference words were available through copy edits, four references appeared among up to 32 masked visual proposals and none among the parent's copy-only proposals. The masked writer ranked none of the four first. Neither exchanged input produced any of the four reference matches in the first 32. The two searches differ in their proposal rules and writable outputs.

The masked visual writer predicted the reference word count on 91 difficult turns, compared with 65 for the causal control. Giving it the correct length on all 435 difficult turns within its capacity yielded just one normalized match. In 393 cases, a word image committed on the first pass already differed from the literal reference. Even when supplied all other human answer words, the writer recovered 677 of 1,937 individually copyable target word images and none of the 327 absent images. On the latter it chose spelling ten times. Forcing the spelling route would recover nine. Among the absent images, 245 had appeared in fitting answers, including 152 among the 2,048 most frequent fitting word images. These counts identify two distinct problems for free writing: locating the right word among visible alternatives and producing a fitting word that the current passage does not contain.
